%============================================================
% BAB V - EVALUASI DAN HASIL (kerangka)
%============================================================
\chapter{EVALUASI DAN HASIL}
\label{chap:evaluasi}

Bab ini melaporkan hasil evaluasi keempat kondisi eksperimen (A, B, C, dan D) pada sampel uji berdasarkan desain evaluasi pada Bab~\ref{chap:analisis-desain} dan implementasi pada Bab~\ref{chap:implementasi}. Seluruh parameter telah dibekukan sebelum evaluasi dijalankan, dan sampel uji hanya dijalankan satu kali. Hasil dilaporkan secara deskriptif, yaitu berupa nilai rata-rata, median, persentase, dan selisih antarkondisi. Pembahasan disusun sebagai berikut.
\begin{enumerate}[1.]
    \item Pelaksanaan evaluasi, yaitu konfigurasi eksperimen pada sampel uji, cakupan data, dan verifikasi kontrol eksperimen.
    \item Keandalan panel \emph{judge} yang digunakan untuk menilai halusinasi, \emph{faithfulness}, dan \emph{citation support}.
    \item Hasil kualitas jawaban, yaitu \emph{cosine similarity}, \emph{hallucination rate}, dan \emph{answer relevance}.
    \item Hasil kualitas \emph{grounding} dan keterlacakan, yaitu \emph{faithfulness}, \emph{citation support}, dan \emph{valid citation rate} (NF2).
    \item Hasil kualitas \emph{retrieval}, yaitu \emph{context relevance} dan kecukupan konteks.
    \item Hasil latensi (NF3).
    \item Analisis perbandingan antarkondisi, yaitu Kondisi C terhadap Kondisi B dan Kondisi D.
    \item Analisis kasus berhasil dan kasus gagal.
    \item Pembahasan, mencakup pemenuhan kriteria sukses NF1, NF2, dan NF3, jawaban atas rumusan masalah, serta posisi hasil terhadap penelitian terdahulu.
    \item Ancaman terhadap validitas.
\end{enumerate}

%============================================================
\section{Pelaksanaan Evaluasi}
\label{sec:pelaksanaan-evaluasi}
%============================================================

\subsection{Konfigurasi Run Resmi}
\label{subsec:konfigurasi-run}
todo...

\subsection{Cakupan Data dan Kelengkapan Run}
\label{subsec:kelengkapan-run}
todo...


\subsection{Verifikasi Kontrol Eksperimen}
\label{subsec:verifikasi-kontrol}
todo...


%============================================================
\section{Keandalan Panel \emph{Judge}}
\label{sec:keandalan-judge}
todo...
%============================================================

%============================================================
\section{Hasil Kualitas Jawaban}
\label{sec:hasil-kualitas-jawaban}
todo...
%============================================================

\subsection{\emph{Cosine Similarity}}
\label{subsec:hasil-cosine}
todo...

\subsection{\emph{Hallucination Rate}}
\label{subsec:hasil-halusinasi}
todo...

\subsection{\emph{Answer Relevance}}
\label{subsec:hasil-answer-relevance}
todo...

%============================================================
\section{Hasil Kualitas \emph{Grounding} dan Keterlacakan}
\label{sec:hasil-grounding}
todo...
%============================================================

\subsection{\emph{Faithfulness}}
\label{subsec:hasil-faithfulness}
todo...


\subsection{\emph{Citation Support}}
\label{subsec:hasil-citation-support}
todo...


\subsection{\emph{Valid Citation Rate} (NF2)}
\label{subsec:hasil-nf2}

todo...

%============================================================
\section{Hasil Kualitas \emph{Retrieval}}
\label{sec:hasil-retrieval}
todo...
%============================================================

\subsection{\emph{Context Relevance} per Item}
\label{subsec:hasil-context-relevance}

todo...

\subsection{Kecukupan Konteks}
\label{subsec:hasil-kecukupan-konteks}
todo...


%============================================================
\section{Hasil Latensi (NF3)}
\label{sec:hasil-latensi}
todo...
%============================================================



%============================================================
\section{Analisis Perbandingan Antar-Kondisi}
\label{sec:analisis-antarkondisi}
todo...
%============================================================

\subsection{Kondisi C terhadap Kondisi B}
\label{subsec:analisis-c-b}
todo...


\subsection{Kondisi C terhadap Kondisi D}
\label{subsec:analisis-c-d}
todo...


\subsection{Ringkasan Lintas Metrik}
\label{subsec:ringkasan-lintas-metrik}
todo...


%============================================================
\section{Analisis Kasus}
\label{sec:analisis-kasus}
todo...
%============================================================

\subsection{Kasus Berhasil}
\label{subsec:kasus-berhasil}

todo...

\subsection{Kasus Gagal dan Pola Kegagalan}
\label{subsec:kasus-gagal}
todo...


%============================================================
\section{Pembahasan}
\label{sec:pembahasan}
%============================================================

\subsection{Pemenuhan Kriteria Sukses Kebutuhan Nonfungsional}
\label{subsec:pembahasan-nonfungsional}
todo...

\subsection{Jawaban atas Rumusan Masalah}
\label{subsec:pembahasan-rumusan-masalah}
todo...


\subsection{Posisi terhadap Penelitian Terdahulu}
\label{subsec:pembahasan-penelitian-terdahulu}
todo...


\subsection{Implikasi dan Klaim Kontribusi}
\label{subsec:pembahasan-implikasi}
todo...


%============================================================
\section{Ancaman terhadap Validitas}
\label{sec:ancaman-validitas}
%============================================================

%============================================================
\section{Ringkasan}
\label{sec:ringkasan-evaluasi}
%============================================================

